% ---- discussion.tex ----
% Data sources: /home/sharaths/projects/PSALM-integration/research/memory/findings.md (F1/F2/F3 contrasts)
% /home/sharaths/projects/PSALM-integration/site/src/data/results.json (BabyLM baseline 74.53, 65.08)
% /home/sharaths/projects/PSALM-integration/data/hf_export/prabhasa_b_ss_0.1_glue/glue_summary.json (GLUE 0.5807)

\section{Discussion}

\subsection{Architecture and Objective as Primary Gains}

The validated findings establish that the largest and most robust improvements in the Prabhasa pipeline come from two sources: the positional encoding (RoPE) and the choice of pretraining objective. RoPE, which applies complex-number rotations in embedding space, provides compositional structure that carries across scales, consistent with its widespread adoption in recent language model families (LLaMA, Qwen). The objective architecture shows a marked scale dependency: at the 10M-parameter regime, the hybrid MLM+CLM objective is statistically indistinguishable from pure MLM (3-seed mean 63.58 \ensuremath{\pm} 1.73 vs.\ 64.09 \ensuremath{\pm} 0.26 BLiMP, overlapping 95\% confidence intervals), while at 100M parameters, pure MLM yields a large gain (73.06 vs.\ 67.57, a +5.49 percentage point difference). This gap reflects a fundamental trade-off: the CLM head's additional gradient signal may stabilize learning on small models where the MLM signal is sparse, but becomes a bottleneck as the MLM objective grows in strength and signal-to-noise. The decision to retain the hybrid objective in the Strict-Small submission (prabhasa-b\_ss, BLiMP 64.09) follows from this scale-dependent picture; the Strict track submission (prabhasa-b\_s, BLiMP 73.06) adopts pure MLM to exploit the larger scale.

\subsection{Interpretation of Null Findings: Scope and Scale Conditions}

Two findings emerge as null or non-significant, and merit careful interpretation to distinguish between genuine refutations of the linguistic hypothesis and results that reflect scope limitations.

\subsubsection{Kāraka Masking: Causal Null at Matched Budget}

The kāraka role-stratified masking scheme (Arm K) yields no measurable causal advantage over uniform random masking (Arm C) when both are held at identical training budget, model size, and corpus. On the targeted 20-paradigm subset (agreement and argument-structure phenomena), the means differ by only +0.10 percentage points (K: 82.03, C: 81.93, 95\% CI [−0.99, +1.20], paired bootstrap over paradigms). Full BLiMP shows a similar pattern (K: 71.77, C: 70.49). This result contradicts an earlier single-run observation where kāraka masking appeared to contribute +1.23 percentage points to overall performance. Root-cause analysis revealed that the earlier figure was confounded: the difference lay not in the role distribution of masking but in two operational details, a higher effective mask rate (freq\_alpha 0.5) and a non-budget-matched masking probability distribution. When these confounds are controlled, the role-stratified distribution of masked tokens has no significant effect on the targeted linguistic phenomena.

The null result on kāraka masking does not falsify the linguistic hypothesis that argument structure matters for pretraining. Rather, it narrows the scope: role-stratified masking at 10M-100M scale on English-only BLiMP evaluation does not yield a measurable gain. The masking scheme may still be beneficial in other contexts: larger models (1B+), longer pretraining runs (100M+ tokens), multilingual settings where cross-lingual transfer brings Sanskrit structure to bear, or evaluation on syntax-specific compositional tasks (SCAN, COGS) that explicitly reward role awareness. A single-seed arm run (Arm K at 100M with RoPE and N-hot embeddings) provides no basis for ranking hypotheses about whether a full power study (e.g., 3-5 seeds) or a larger-scale confirmation would reverse the null. The pre-registered, budget-matched comparison, however, provides a valid and conservative estimate under the stated conditions: role-stratified masking does not move the needle on BLiMP agreement and argument-structure tasks at matched budget.

\subsubsection{Auxiliary Objective: Weak and Non-Significant at 5 Seeds}

A supervised auxiliary objective that predicts kāraka roles at the token level showed a promising single-seed result ($+2.65$ percentage points, 95\% CI [+1.13, +4.19]) on the targeted subset. Extension to 5 pre-registered seeds (seed0: +2.65, seed1: +0.61, seed2: +1.12, seed3: −0.22, seed4: −0.34) revealed substantial inter-seed variance. The 5-seed mean difference was +0.76 (95\% CI [−0.74, +2.27], $t = 1.41$, $\text{df} = 4$, $p > 0.05$), not statistically significant. The effect attenuated monotonically with sample size: +2.65 at 1 seed, +1.46 at 3 seeds (Table in findings.md), +0.76 at 5 seeds, suggesting seed0 represented a positive fluctuation rather than a robust phenomenon.

A specificity control using shuffled-role labels revealed that the role-label signal is learnable (the model achieves lower training loss with genuine role annotations than with random labels), yet the shuffled-role auxiliary provides equivalent downstream benefit (+0.58 on the targeted subset, vs.\ +0.76 for real roles), indicating the effect is not specifically tied to linguistic role structure. The auxiliary mechanism thus appears to provide weak and nonspecific regularization, not a causal improvement from role-aware supervision.

These null results can be interpreted as scope limitations. The 10M-scale regime, while the published BabyLM budget, may lack sufficient corpus breadth and model capacity for the auxiliary objective to recover compositional benefits. The F1 scale-dependence finding suggests that objective effects can interact with model size; a 100M-scale replication of the auxiliary objective, though computationally costly, would determine whether the null is fundamental or scale-dependent. Alternatively, evaluation surfaces beyond BLiMP may be more sensitive: the targeted agreement and argument-structure paradigms are narrow, and a broader evaluation on SCAN or COGS (which directly test compositional generalization to novel argument combinations) might reveal benefits that BLiMP does not capture.

\subsection{Integration with Prior BabyLM Findings}

The Prabhasa results extend and refine prior work on data-efficient pretraining. Warstadt et al.~\cite{Warstadt2023} and the BabyLM shared task~\cite{Hu2024BabyLM} have established that objective choice, corpus composition, and optimizer selection are high-leverage parameters at small scales. This work confirms that objective design is scale-sensitive~\cite{Huebner2021} and~\cite{CharpentierSamuel2024}, and that positional encodings (RoPE) deliver consistent gains across configurations. The null findings on kāraka masking align with prior observations that structure-aware masking provides modest and context-dependent benefits: Charpentier and Samuel~\cite{CharpentierSamuel2023} on morphologically-aware pretraining show gains, but primarily when combined with explicit morphological vocabularies or linguistically-rich evaluation sets; Clark et al.~\cite{Clark2020} on ELECTRA-style discrimination similarly find that structure must be paired with sufficient model scale or evaluation breadth to manifest. Prabhasa's results suggest that in the 10M English-only microworld, such structure provides interpretability and modularity without immediate downstream lift on standard diagnostics, a finding consistent with the broader trend that small models saturate quickly on BLiMP and require scale or task specificity to benefit from structural priors.

\subsection{Methodological Contributions: Pre-Registration and Matched-Budget Controls}

Beyond the empirical findings, the Prabhasa pipeline demonstrates a methodological contribution to rigorous pretraining research. Pre-registration of hypotheses and effect-size thresholds (F1: scale-dependent objective; F2: causality at matched budget; F3: 5-seed auxiliary objective) enforces discipline and prevents outcome-driven narrative construction. The matched-budget comparisons (Arm K and Arm C in F2 held at identical mask volume, parameter count, and corpus) eliminate confounds that could otherwise inflate estimates of mechanism efficacy. The extension of F3 from 1 seed to 5 seeds, coupled with specificity controls (shuffled-role negative control), exemplifies the multi-seed regimen recommended by the Synthetic Data and Fine-Tuning literature and the BabyLM challenge itself. This approach prevents cherry-picking of favorable seeds and exposes inter-seed variance that single-seed publications often mask.

The decision to report a null result with full transparency (seed-level values, confidence intervals, specificity controls) rather than to frame the null as ``no evidence of effect'' or to emphasize exploratory directions contributes to a research culture in which null findings are validated research outcomes, not failures. In the small-model regime, where compute is constrained and replication is costly, clarity about what does and does not work is valuable for downstream researchers.

\subsection{Threats to Validity}

Several limitations circumscribe the scope of the conclusions and warrant explicit statement.

\emph{Scale of evaluation and generalization.} The BabyLM fixed budget (10M tokens, 100M parameters) is a constrained microworld, orders of magnitude smaller than contemporary pretraining runs. Whether the null findings on kāraka masking and auxiliary objectives generalize to 1B-scale models or 100M-token pretraining runs is unknown. The scale-dependent objective effect (F1) suggests that mechanism efficacy may shift with model capacity; future 100M-scale evaluation of the auxiliary objective would address this question, but the computational cost (estimated at 2-3 GPU-days on DGX hardware for each configuration) is deferred to future phases.

\emph{Kāraka role accuracy.} The kāraka role assignments used in the masking and auxiliary objectives are derived from BPE-token-level heuristics: token identity matching to a small Sanskrit lexicon and frequency-based role assignment. This contrasts with full syntactic parsing (which would employ spaCy or Universal Dependencies on the Sanskrit corpus), which achieves higher accuracy on role recovery. The current heuristic approach introduces noise into the role signal; a robustness check using high-accuracy parses (via Vidyut full morphological analysis on the Sanskrit corpus) would determine whether the F2 and F3 nulls partly reflect role-assignment noise or are genuine.

\emph{Auxiliary objective design.} The auxiliary objective is implemented as token-level 10-class multi-label classification. Alternative designs (continuous role embeddings, hierarchical classification, phrasal-level role prediction) may exhibit different behavior. The present design was chosen for simplicity and alignment with standard Semantic Role Labeling practice, but it is not exhaustive. A design sensitivity analysis is left to future work.

\emph{Seed variance and statistical power.} The F3 5-seed study is powered to detect main effects of approximately +2 percentage points or larger ($\alpha = 0.05$, $n = 5$). The observed seed variance (ranging from −0.34 to +2.65) is substantial, indicating that smaller effect sizes (e.g., +0.5--1.0 pp) cannot be reliably resolved with 5 seeds. A higher-powered replication (10--15 seeds) would better characterize the true effect distribution and narrow the confidence intervals.

\emph{Single-seed Strict-track run.} The primary Strict-track model (prabhasa-b\_s, BLiMP 73.06 at 100M) is a single seed. Computational constraints prevented 3-5 seed replication at the larger scale. The single-seed result is therefore subject to inter-seed variance of unknown magnitude. A ceiling estimate can be drawn from the Strict-Small 10M runs, where 3-seed standard deviations range from 0.26 to 1.73; a similar variance at 100M would yield confidence intervals of approximately \ensuremath{\pm}1.5 percentage points. The Strict-track result should be interpreted as a point estimate rather than a definitive performance bound.

\emph{Proxy evaluation scale.} The dose-arm comparison (Arms A--D in the pre-pretraining phase) was conducted at 60M parameter scale, not at the target 100M or larger. The null signal at proxy scale (all arms within 0.6 pp) motivated the decision not to pursue a 350M confirmation. This is a reasonable heuristic but not definitive: the null may reflect saturation of the 60M regime on BLiMP rather than a genuine null at larger scales.

\subsection{Implications for Research Direction}

The Prabhasa findings suggest several productive avenues for future work. First, the scale-dependent objective effect (F1) motivates a systematic characterization of the objective-scale interaction: at what model size does pure MLM become superior, and does the crossover differ by evaluation task (BLiMP vs.\ SCAN vs.\ COGS)? Second, the null on kāraka masking at matched budget does not rule out role-aware structure as beneficial; it suggests that the benefit (if any) may be task-specific or scale-dependent. Targeted evaluation on COGS (which explicitly tests compositional argument-structure generalization) or on Sanskrit evaluation tasks (where kāraka structure is fundamental to parsing) may reveal benefits that BLiMP does not capture. Third, the weak auxiliary objective result may reflect the choice of supervision signal; alternatives such as predicting higher-level semantic frames or using contrastive objectives (e.g., InfoNCE over role-aligned pairs) may be more effective. Fourth, multilingual pretraining with Sanskrit as a bridge language may recover structural benefits that are lost in English-only training.

Finally, the methodological contribution of pre-registration and matched-budget controls deserves replication in other pretraining studies. The small-model regime, where compute is constrained and publications often rely on single-seed runs, is particularly vulnerable to selective reporting and confounding. Adopting pre-registration, matched-budget controls, and transparent reporting of null findings would improve the credibility and reproducibility of the BabyLM literature.

\subsection{Conclusion}

The Prabhasa framework demonstrates that robust gains in sample-efficient pretraining come from architecture (RoPE) and objective design (pure-MLM at scale), while linguistically-motivated mechanisms for masking and supervision deliver interpretability and design coherence without measured downstream lift on standard diagnostics at 10M-100M scale. This result does not dismiss linguistic structure as irrelevant to pretraining; rather, it underscores that the specific mechanisms tested (role-stratified masking, token-level role supervision) require larger models, longer training runs, or task-specific evaluation to manifest empirical benefit. The value of the work lies both in the validated findings (F1 scale-dependent objectives) and in the transparent null results (F2, F3), grounded in pre-registration and matched-budget methodology that the field should adopt more broadly.
